\documentclass[10pt,a4paper]{article}
\usepackage[T1]{fontenc}
\usepackage{lmodern}
\usepackage{amsmath,amssymb,amsthm,mathtools}
\usepackage[margin=24mm]{geometry}
\usepackage{microtype}
\usepackage{booktabs,array,longtable}
\usepackage{xcolor}
\usepackage{hyperref}
\usepackage{fancyhdr}
\hypersetup{colorlinks=true,linkcolor=black,citecolor=black,urlcolor=blue,pdfauthor={Matthew J. Colbrook}}
\pagestyle{fancy}
\fancyhf{}
\fancyhead[L]{\small Partial result / MI-09}
\fancyhead[R]{\small 11 September 2026}
\fancyfoot[C]{\thepage}
\setlength{\headheight}{14pt}
\setlength{\parskip}{4pt}
\setlength{\parindent}{0pt}
\newtheorem{theorem}{Theorem}[section]
\newtheorem{proposition}[theorem]{Proposition}
\newtheorem{lemma}[theorem]{Lemma}
\newtheorem{corollary}[theorem]{Corollary}
\theoremstyle{remark}
\newtheorem{remark}[theorem]{Remark}
\newcommand{\M}{\mathbb M}
\newcommand{\C}{\mathbb C}
\newcommand{\R}{\mathbb R}
\newcommand{\tr}{\operatorname{tr}}
\newcommand{\diag}{\operatorname{diag}}
\newcommand{\spec}{\operatorname{spec}}
\newcommand{\rank}{\operatorname{rank}}
\newcommand{\abs}[1]{\lvert #1\rvert}
\newcommand{\norm}[1]{\lVert #1\rVert}
\newcommand{\opnorm}[1]{\lVert #1\rVert_{\infty}}
\newcommand{\schatten}[2]{\lVert #1\rVert_{#2}}
\newcommand{\symmod}{\mathcal S}
\newcommand{\maxmod}{\mathcal M}
\newcommand{\draftnotice}{\begin{center}\small\textit{Proposed mathematical resolution. Complete proof supplied; not submitted or peer reviewed.}\end{center}}

\usepackage{xurl}
\setlength{\emergencystretch}{3em}
\allowdisplaybreaks[1]

\makeatletter
\renewcommand{\@maketitle}{\newpage\null\begin{center}
{\Large\@title\par}\vskip .8em
{\normalsize\begin{tabular}[t]{c}\@author\end{tabular}\par}\vskip .5em
{\small\@date}\end{center}\par\vskip .5em}
\makeatother

\title{MI-09: Partial Result}
\author{Matthew J. Colbrook\\
\small Department of Applied Mathematics and Theoretical Physics\\
\small University of Cambridge, Cambridge, United Kingdom\\
\small\href{mailto:m.colbrook@damtp.cam.ac.uk}{m.colbrook@damtp.cam.ac.uk}}
\date{11 September 2026}
\begin{document}
\maketitle
\textbf{Repository status: Partially resolved.}
The complete argument received an independent Codex-agent PASS review
for the stated partial result only.
The original draft was AI-assisted; the reviewed proof below is unchanged.
This is independent agent verification, not external human peer review or
formal proof certification. \href{https://github.com/ajt60gaibb/OpenProblemsInNLA/blob/main/references/colbrook-matrix-2026-09-11/verification/reviews/MI-09-review.md}{Detailed review and scope}.
\par\medskip
% BEGIN REVIEWED PROOF
\section{MI-09, dimension-two subproblem: an exact sharp universal norm constant}
\label{sec:mi09}
Write $\symmod(T)=(|T|+|T^*|)/2$. We determine the sharp constant uniform over all unitarily invariant norms in dimension two. This resolves the dimension-two question in \cite[Problem 1.11]{Zhang2026}, including its operator-norm case, but does not determine the separate optimal constants for every Schatten exponent and every matrix dimension in the broader MI-09 program.

\begin{theorem}\label{thm:mi09}
For every $m\geq2$, every $T_1,\ldots,T_m\in\M_2(\C)$, and every unitarily invariant norm $\|\!\|\!|\cdot|\!\|\!\|$,
\[
 \|\!\|\!|\symmod(\textstyle\sum_jT_j)|\!\|\!\|
 \leq \sqrt{6\sqrt3-9}\,
 \|\!\|\!|\textstyle\sum_j\symmod(T_j)|\!\|\!\|.
\]
The constant $\sqrt{6\sqrt3-9}=1.1799596795\ldots$ is optimal when the inequality is required for all such norms. It is already attained in the operator norm by two real rank-one matrices. In fact, the sharp numerical-radius inequality is
\begin{equation}\label{eq:mi09-radius}
 \opnorm{\symmod(T)}\leq\sqrt{6\sqrt3-9}\,w(T),\qquad T\in\M_2(\C),
\end{equation}
where $w(T)=\max_{\|v\|=1}|v^*Tv|$.
\end{theorem}

\subsection{Reduction to a numerical-radius inequality}
For every matrix $Z$ and unit vector $v$, the polar decomposition and Cauchy--Schwarz give
\[
 |v^*Zv|\leq\sqrt{(v^*|Z|v)(v^*|Z^*|v)}\leq v^*\symmod(Z)v.
\]
Consequently, with $T=\sum_jT_j$ and $D=\sum_j\symmod(T_j)$,
\[
 w(T)\leq\opnorm{D}.
\]
Thus \eqref{eq:mi09-radius} implies the claimed operator-norm inequality. Also
\[
 \tr\symmod(T)=\|T\|_1\leq\sum_j\|T_j\|_1=\tr D.
\]
The largest-eigenvalue inequality and this trace inequality imply weak majorization of the two eigenvalues of $\symmod(T)$ by those of $\sqrt{6\sqrt3-9}\,D$. Hence the inequality holds for every unitarily invariant norm.

\subsection{A normal form adapted to the symmetric modulus}
For a nonzero $2\times2$ matrix, set $q=\|T\|_1/2$. The elementary two-dimensional square-root identity yields
\begin{equation}\label{eq:mi09-mod-formula}
 \symmod(T)=\frac{T^*T+TT^*+2|\det T|I}{2\|T\|_1}.
\end{equation}
Indeed, for a positive $2\times2$ matrix $P$, its positive square root equals
$(P+\sqrt{\det P}\,I)/\sqrt{\tr P+2\sqrt{\det P}}$; the formula extends to singular matrices by continuity.

Multiplication by a scalar of modulus one preserves both sides of \eqref{eq:mi09-radius}, as does unitary similarity. Make $\tr T=2h$ real and positive, and diagonalize $\symmod(T)$ in decreasing order. The off-diagonal entry in the numerator of \eqref{eq:mi09-mod-formula} is $2h(t_{12}+\overline{t_{21}})$, so $t_{21}=-\overline{t_{12}}$. A diagonal unitary similarity now gives
\begin{equation}\label{eq:mi09-normal}
 T=\begin{pmatrix}h+x+ik&b\\-b&h-x-ik\end{pmatrix},
 \qquad h,x,b\geq0,\quad k\in\R.
\end{equation}
The ordering of the symmetric-modulus eigenvalues gives $x\geq0$. Cases with zero trace or a repeated symmetric-modulus eigenvalue follow by continuity from these generic cases.

Let $u=h^2+x^2+k^2+b^2$. Direct calculation gives
\[
 q^2=\frac{u+|\det T|}{2},\qquad
 \symmod(T)=\diag\left(q+\frac{hx}{q},q-\frac{hx}{q}\right).
\]
Moreover,
\[
 (q^2-x^2)(q^2-h^2-b^2)=q^2k^2.
\]
We have $h\leq q$ by $|\tr T|\leq\|T\|_1$, and $x\leq q$ follows from the displayed formula for $q^2$ (or by evaluating its quadratic equation at $x^2$). By homogeneity take $q=1$. By continuity we may assume $0<h,x<1$. Then
\begin{equation}\label{eq:mi09-b}
 b^2=1-h^2-\frac{k^2}{1-x^2},\qquad
 \opnorm{\symmod(T)}=1+hx.
\end{equation}
For the normal form \eqref{eq:mi09-normal}, the numerical range is the filled ellipse
\[
 \{h+xs+i(ks+bt):s^2+t^2\leq1\}.
\]
This follows directly by writing the expectation of the real and imaginary Hermitian parts on a unit vector in $\C^2$.

\subsection{A sharp lower bound for the numerical radius}
Put
\[
 H=h^2,\quad X=x^2,\quad P=1-X,\quad Q=1-H,\quad
 g=1-H-X,\quad z=k^2.
\]
First suppose $g\geq hx$. We claim
\begin{equation}\label{eq:mi09-wlower}
 w(T)^2\geq\frac{(1-H)(1-X)}{1-H-X}
 =1+\frac{HX}{g}.
\end{equation}
The generic case has $g>0$. To prove the claim, set
\[
 M=\begin{pmatrix}X+z&kb\\kb&b^2\end{pmatrix},\qquad a=\begin{pmatrix}hx\\0\end{pmatrix}.
\]
The ellipse formula gives the elementary quadratic maximization identity
\begin{equation}\label{eq:mi09-trust}
 w(T)^2=H+\inf_{\lambda>\lambda_{\max}(M)}
 \left\{\lambda+a^T(\lambda I-M)^{-1}a\right\}.
\end{equation}
Here the infimum may be attained only as a boundary limit. For completeness, completion of the square gives the upper bound in \eqref{eq:mi09-trust} for every $\lambda>\lambda_{\max}(M)$. Equality holds when $v=(\lambda I-M)^{-1}a$ has norm one. Its norm decreases continuously to zero as $\lambda\to\infty$. If it does not reach one before the lower endpoint, its endpoint value has norm at most one and is orthogonal to a top eigenvector; adding a suitable multiple of that eigenvector gives a unit vector achieving equality in the endpoint limit. This proves \eqref{eq:mi09-trust}, including that boundary case.

Equation \eqref{eq:mi09-b} gives
\[
 \det(QI-M)=-\frac{HXz}{P}\leq0.
\]
Since $Q>X$ here, it follows that $\lambda_{\max}(M)\geq Q$. Write $\lambda=Q+d$ with $d>0$, and put
\[
 \Delta=\det(\lambda I-M)=d(g+d)+\frac{zX(d-H)}P>0.
\]
The expression inside the infimum, including $H$, becomes
\[
 f=1+d+\frac{HX(d+z/P)}{\Delta}.
\]
The following exact polynomial identity is the key estimate:
\begin{align*}
 g\Delta\left(f-1-\frac{HX}{g}\right)
 ={}&d^2(g^2-HX+gd)\\
 &+\frac{zX}{P}\bigl(gd^2-HQd+HQP\bigr).
\end{align*}
The first term is nonnegative because $g\geq hx$. Also
$HQ\leq gP$, since $gP-HQ=g^2-HX\geq0$.
If $0\leq d\leq P$, then
$gd^2-HQd+HQP=gd^2+HQ(P-d)\geq0$.
If $d\geq P$, then
$gd^2-HQd+HQP=d(gd-HQ)+HQP\geq0$.
Thus $f\geq1+HX/g$ for every admissible $\lambda$. Taking the infimum proves \eqref{eq:mi09-wlower}.

Let $p=hx$. The condition $g\geq p$ implies $0\leq p\leq1/3$. Since $H+X\geq2p$, equation \eqref{eq:mi09-wlower} yields
\[
 \frac{\opnorm{\symmod(T)}^2}{w(T)^2}
 \leq \frac{(1+p)^2g}{g+p^2}
 \leq\frac{(1+p)^2(1-2p)}{(1-p)^2}.
\]
The last rational function on $[0,1/3]$ has its maximum at $p=2-\sqrt3$, where its value is $6\sqrt3-9$. This follows by differentiating: the only critical point in that interval solves $p^2-4p+1=0$.

It remains to consider $g\leq hx$. Now $w(T)\geq h+x$, using the real part of $T$. If $p=hx\leq1/3$, then
\[
 (h+x)^2=H+X+2p\geq1+p,
 \qquad\frac{(1+p)^2}{(h+x)^2}\leq1+p\leq\frac43.
\]
If $p\geq1/3$, then
\[
 \frac{(1+p)^2}{(h+x)^2}\leq\frac{(1+p)^2}{4p}\leq\frac43
 \qquad(1/3\leq p\leq1).
\]
Since $4/3<6\sqrt3-9$, this proves \eqref{eq:mi09-radius} in all cases.

\subsection{Two explicit rank-one summands attaining the bound}
Let
\[
 p=2-\sqrt3,\qquad h=\sqrt p,\qquad b=\sqrt{1-p},\qquad
 w=\frac{1-p}{\sqrt{1-2p}},\qquad z=\frac{h\sqrt{1-p}}{\sqrt{1-2p}}.
\]
Define the real matrices
\[
 T_1=\frac12\begin{pmatrix}w+2h&b+z\\-b-z&-w\end{pmatrix},\qquad
 T_2=\frac12\begin{pmatrix}2h-w&b-z\\-b+z&w\end{pmatrix}.
\]
The identities $b^2+z^2=w^2$ and $bz=wh$ show that both determinants vanish. Directly, their nonzero singular values are $w+h$ and $w-h$, respectively. Applying \eqref{eq:mi09-mod-formula} gives
\[
 \symmod(T_1)=\frac{w+h}{2}I+\frac h2\diag(1,-1),\qquad
 \symmod(T_2)=\frac{w-h}{2}I-\frac h2\diag(1,-1).
\]
Their sum is $wI$, whereas
\[
 T_1+T_2=\begin{pmatrix}2h&b\\-b&0\end{pmatrix},\qquad
 \symmod(T_1+T_2)=\diag(1+p,1-p).
\]
Therefore
\[
 \frac{\opnorm{\symmod(T_1+T_2)}}{\opnorm{\symmod(T_1)+\symmod(T_2)}}
 =\frac{1+p}{w}=\sqrt{6\sqrt3-9}.
\]
Padding with zero summands handles every $m\geq2$. This proves sharpness and completes the proof of Theorem~\ref{thm:mi09}.\hfill$\square$

% END REVIEWED PROOF
\begin{thebibliography}{1}

\bibitem{Zhang2026}
Teng Zhang.
\newblock Operator symmetric moduli and sharp triangle inequalities.
\newblock arXiv:2603.01046, 2026.
\newblock Version consulted: v1; Conjecture 1.6, equation (1.4), Problems 1.11
  and 1.17.
\newblock URL: \url{https://arxiv.org/abs/2603.01046}.

\end{thebibliography}

\end{document}
